@@M0@@
\item
\emph{@@M1@@ is numerically equivalent to @@M2@@.}
\end{enumerate}
\emph{If @@M3@@ is projective and equipped with an ample invertible module @@M4@@, the preceding conditions are also equivalent to the following}:
\begin{enumerate}[label=(\roman*),start=4]
\item \emph{For every integer @@M5@@ (positive or negative), @@M6@@ is ample.}
\item \emph{(If @@M7@@ is irreducible). For every pair of integers @@M8@@, we have}
@@M9@@
\item \emph{(If @@M10@@ is irreducible of dimension @@M11@@). We have}
@@M12@@
\end{enumerate}

Let us prove that (i) implies (ii). Applying the d\'evissage lemma (EGA III 3.1) and 4.5 (i), we reduce to the case where @@M13@@ is integral and @@M14@@. We must then prove that the polynomial @@M15@@ is constant. For this, we may replace @@M16@@ by a multiple, and thus assume that @@M17@@ is algebraically equivalent to @@M18@@. In this case, if @@M19@@ is deformed into @@M20@@ by the algebraic family defined by an invertible sheaf @@M21@@ on some @@M22@@ where @@M23@@ is connected, of finite type, the polynomial @@M24@@ does not depend on @@M25@@ (EGA III 7.9.5); hence the implication.

Let @@M26@@ be a closed integral curve in @@M27@@. Taking @@M28@@ in (ii), we find @@M29@@. Thus the formula (``Riemann's theorem'') @@M30@@ gives @@M31@@; hence (ii) implies (iii).

Let @@M32@@ be a Chow presentation: @@M33@@ is projective and @@M34@@ is surjective. If we grant that (iii) implies (i) for @@M35@@, then
